% Generado por scripts/tablas.py a partir de Practica3_pendulos_acoplados.xlsx. No editar a mano.

\begin{table}[ht]
\centering
\caption{Periodos y frecuencias angulares para cada longitud de acoplamiento. $T_1$ es la media de los dos péndulos en contrafase; $\omega_1=2\pi/T_1$, $\omega_{-}=\pi/T_{\mathrm{b}}$ y $\omega_{+}=2\pi/T_{+}$. Las incertidumbres se tratan en la \cref{sec:ajustes}.}
\label{tab:frecuencias}
\renewcommand{\arraystretch}{1.22}\setlength{\tabcolsep}{5pt}
\begin{tabular}{@{}l c c c c c c c c@{}}
\toprule
$\ell$ & $\ell^{2}$ (\si{\meter\squared}) & $T_1$ (\si{\second}) & $\omega_1$ (\si{\radian\per\second}) & $\omega_1^{2}$ (\si{\radian\squared\per\second\squared}) & $T_{\mathrm{b}}$ (\si{\second}) & $\omega_{-}$ (\si{\radian\per\second}) & $T_{+}$ (\si{\second}) & $\omega_{+}$ (\si{\radian\per\second}) \\
\midrule
$\ell_{1}$ & \num{0.0625} & \num{1.9620} & \num{3.2024} & \num{10.256} & \num{105.09} & \num{0.0299} & \num{1.9736} & \num{3.1836} \\
$\ell_{2}$ & \num{0.1225} & \num{1.9288} & \num{3.2576} & \num{10.612} & \num{56.09} & \num{0.0560} & \num{1.9560} & \num{3.2123} \\
$\ell_{3}$ & \num{0.2025} & \num{1.8912} & \num{3.3223} & \num{11.038} & \num{35.46} & \num{0.0886} & \num{1.9312} & \num{3.2535} \\
$\ell_{4}$ & \num{0.3025} & \num{1.8820} & \num{3.3386} & \num{11.146} & \num{26.12} & \num{0.1203} & \num{1.9232} & \num{3.2670} \\
$\ell_{5}$ & \num{0.4225} & \num{1.8452} & \num{3.4052} & \num{11.595} & \num{17.80} & \num{0.1765} & \num{1.9080} & \num{3.2931} \\
\bottomrule
\end{tabular}
\end{table}
